% ==============================================================================
% capitulos/03-tecnologias.tex - Infraestructura y Tecnologías
% ==============================================================================

\chapter{Infraestructura y Tecnologías}
\label{cap:tecnologias}

\section{Estudio de Arquitecturas de Software}
\label{sec:estudio_arquitecturas}

\subsection{Introducción a las Arquitecturas de Software}
\begin{itemize}
    \item Definición de arquitectura de software
    \item Importancia en el desarrollo de sistemas
    \item Factores de calidad:
    \begin{itemize}
        \item Rendimiento
        \item Escalabilidad
        \item Mantenibilidad
    \end{itemize}
    \item Clasificación general de arquitecturas:
    \begin{itemize}
        \item Monolíticas
        \item Distribuidas
        \item Orientadas a servicios
        \item Basadas en eventos
    \end{itemize}
\end{itemize}

\subsection{Arquitectura Monolítica}
\begin{itemize}
    \item Definición y estructura
    \item Funcionamiento interno
    \item Ventajas
    \item Desventajas
    \item Casos de uso
\end{itemize}

\subsection{Arquitectura Orientada a Servicios (SOA)}
\begin{itemize}
    \item Definición y principios
    \item Enterprise Service Bus (ESB)
    \item Protocolos de comunicación:
    \begin{itemize}
        \item SOAP
        \item XML
    \end{itemize}
    \item Ventajas
    \item Desventajas
    \item Diferencias con microservicios
\end{itemize}

\subsection{Arquitectura de Microservicios}
\begin{itemize}
    \item Definición y características
    \item Principios fundamentales:
    \begin{itemize}
        \item Independencia de despliegue
        \item Alta cohesión
        \item Bajo acoplamiento
    \end{itemize}
    \item Descomposición funcional
    \item Comunicación entre servicios:
    \begin{itemize}
        \item REST
        \item JSON
        \item API Gateway
    \end{itemize}
    \item Patrones clave:
    \begin{itemize}
        \item API Gateway
        \item Service Discovery
        \item Circuit Breaker
    \end{itemize}
    \item Gestión de datos:
    \begin{itemize}
        \item Base de datos por servicio
        \item Base de datos compartida
    \end{itemize}
    \item Ventajas
    \item Desventajas
\end{itemize}

\subsection{Arquitectura de Microservicios en 3 Capas}
\begin{itemize}
    \item Capa de presentación
    \item Capa de lógica
    \item Capa de datos
    \item Flujo de petición
    \item Justificación del diseño
\end{itemize}

\subsection{Arquitectura Dirigida por Eventos (Event-Driven)}
\begin{itemize}
    \item Concepto de eventos
    \item Productores y consumidores
    \item Message brokers:
    \begin{itemize}
        \item RabbitMQ
        \item Kafka
    \end{itemize}
    \item Comunicación asíncrona
    \item Consistencia eventual
    \item Ventajas
    \item Desventajas
\end{itemize}

\subsection{Arquitectura Serverless (FaaS)}
\begin{itemize}
    \item Definición
    \item Ejecución basada en eventos
    \item Plataformas:
    \begin{itemize}
        \item AWS Lambda
        \item Google Cloud Functions
    \end{itemize}
    \item Características:
    \begin{itemize}
        \item Escalabilidad automática
        \item Pago por uso
        \item Stateless
    \end{itemize}
    \item Ventajas
    \item Desventajas
\end{itemize}

\subsection{Arquitectura Hexagonal (Ports and Adapters)}
\begin{itemize}
    \item Definición
    \item Separación del dominio
    \item Puertos y adaptadores
    \item Independencia tecnológica
    \item Beneficios en pruebas
\end{itemize}

\subsection{Comparación de Arquitecturas}
\begin{itemize}
    \item Escalabilidad
    \item Complejidad
    \item Acoplamiento
    \item Mantenibilidad
    \item Casos de uso
\end{itemize}

\subsection{Selección de Arquitectura}
\begin{itemize}
    \item Criterios de selección
    \item Evaluación según contexto
    \item Justificación de la arquitectura elegida
\end{itemize}